The \textsc{bisect-apportion} crate implements Webster's method via the
same priority-queue architecture as Huntington-Hill, with the priority
function changed to:
\[
  P_i^{\text{Web}}(s) = \frac{p_i}{s+0.5}.
\]
The implementation computes this as:

\begin{verbatim}
RoundingRule::Webster => p / (n + 0.5),
\end{verbatim}

where \texttt{p} and \texttt{n} are the \texttt{f64} casts of population
and current seat count.

\subsection{Current API Usage}

\begin{verbatim}
let seats = apportionment_divisor(&populations, 435, RoundingRule::Webster);
\end{verbatim}

There is not yet a standalone \texttt{bisect apportion --method webster} CLI.

\subsection{Test Coverage}

\begin{itemize}
  \item Equal-population and 3:1 ratio synthetic examples cover the Webster path.
  \item Shared divisor-method tests assert total-seat preservation and the
    one-seat constitutional floor.
  \item The paradox module checks that Webster has no Alabama-paradox event over
    the synthetic house-size range 5 through 50.
  \item The package fixture
    \texttt{docs/examples/j-apportionment-evidence-packages/divisor-method-smoke}
    hash-binds a three-state Webster allocation and replays it through
    \texttt{bisect-apportion::divisor\_evidence}.
\end{itemize}

\subsection{Sainte-Lagu\"e Mode}

The \textsc{bisect-apportion} Webster implementation can also be used as a
Sainte-Lagu\"e-style party-list allocation helper when the \texttt{populations}
argument is interpreted as party vote totals and \texttt{house\_size} as the
number of seats to allocate.  This is a library use case, not yet a separate CLI
mode.
